\documentclass[11pt]{article}
\usepackage{amsmath,amssymb,amsthm}
\usepackage[margin=1in]{geometry}
\newtheorem{proposition}{Proposition}
\begin{document}
\noindent\textbf{Official statement (Cell 3).}
\emph{Any $d+1$ lines in $\mathbb R^d$ satisfy
$S\leq\bigl(\binom{d+1}{2}-1\bigr)\pi/2$, i.e. $A\geq\pi/2$.}

\bigskip
\noindent Here $S$ is the sum of the acute angles between all pairs
of lines.  For unit representatives $x_1,\ldots,x_N$, write
\[
 A=\sum_{i<j}\arcsin|\langle x_i,x_j\rangle|.
\]
Since $S=\binom N2\pi/2-A$, the official statement is equivalent to
$A\geq\pi/2$ when $N=d+1$.

\begin{proposition}[Orthogonality saturation]
For fixed $N$ and $d\geq2$, there is a maximizer of $S$ such that
each line is either orthogonal to every other line, or the lines
orthogonal to it span a subspace of dimension at least $d-1$.
\end{proposition}
\begin{proof}
The space of $N$-tuples of lines is a product of compact projective
spheres, and $S$ is continuous.  Choose a maximizer with as many
orthogonal pairs as possible.  Fix its line $\ell_i$ and let $W$ be
the span of the representatives of all lines orthogonal to $\ell_i$.
Suppose $\dim W\leq d-2$ and $\ell_i$ is not orthogonal to every other
line.  A unit representative $x_i$ lies in $W^\perp$, whose dimension
is at least two.  Choose a two-dimensional plane $U\subset W^\perp$
containing $x_i$ such that the inner product with one nonorthogonal
line varies along its unit circle.  Such a plane exists: the
projection of that line onto $W^\perp$ is nonzero, and one can choose
the other direction of $U$ accordingly.  Rotate only $x_i$ in $U$.

Consider the closed interval between the two consecutive positions
where a new orthogonality occurs, one on either side of the original
position.  There are such positions because the varying nonzero
inner product is sinusoidal and has zeros.  Throughout the open
interval, every previously nonorthogonal inner product has fixed
sign.  After a possible sign change, its corresponding angle term
has the form
\[
 f(t)=\arccos\bigl(r\cos(t-\phi)\bigr),\qquad
 0\leq r\leq1,\quad \cos(t-\phi)>0.
\]
For $0<r<1$, direct differentiation gives
\[
 f''(t)=\frac{r(1-r^2)\cos(t-\phi)}
 {(1-r^2\cos^2(t-\phi))^{3/2}}\geq0.
\]
For $r=0$ the term is constant; for $r=1$ it is affine away from
coincidence, with a convex corner at coincidence.  Thus the sum of
all varying angle terms is convex on this interval, including any
coincidence points.  At least one endpoint has value at least the
value at the original position.  Global maximality makes that
endpoint another maximizer.  Every old orthogonality involving
$\ell_i$ persists because the rotation stays in $W^\perp$, while
an additional orthogonal pair appears.  This contradicts our choice
of maximizer.  Hence the stated alternative holds for every line.
\end{proof}

\begin{proposition}[Cell 3]
For every $d\geq1$, any $d+1$ lines in $\mathbb R^d$ satisfy
\[
 A\geq\frac\pi2,
 \qquad
 S\leq\left(\binom{d+1}{2}-1\right)\frac\pi2.
\]
\end{proposition}
\begin{proof}
For $d=1$, the two lines coincide, so $A=\pi/2$.  Suppose $d\geq2$.
It is enough to prove the inequality at a maximum of $S$; choose
the saturated maximizer furnished by the preceding proposition.
Let $J$ join two indices exactly when their lines are
\emph{not} orthogonal.  A vertex adjacent in $J$ to nobody is
harmless.  At every other vertex, the orthogonality neighbors span
dimension at least $d-1$, so there are at least $d-1$ of them among
the other $d$ vertices.  Consequently $J$ has maximum degree one:
it is a matching, together with isolated vertices.

Distinct connected components of $J$ are mutually orthogonal.
Each component consists of one line or a pair of nonorthogonal lines;
it has rank one or two, with a pair having rank one precisely when
the two lines coincide.  If every component had rank equal to its
number of lines, their mutually orthogonal spans would have total
rank $d+1$, impossible in $\mathbb R^d$.  Therefore one paired
component has rank one.  Its two coincident lines contribute
$\arcsin1=\pi/2$ to $A$.  All other contributions are nonnegative,
giving the claim.
\end{proof}

\paragraph{Sharpness and equality at a saturated maximizer.}
Take an orthonormal basis of $\mathbb R^d$ and repeat one basis line.
There is exactly one nonorthogonal pair, which contributes $\pi/2$;
thus equality holds.  Conversely, the matching argument shows that
every \emph{saturated maximizer} with equality has exactly this form,
up to an orthogonal transformation and relabeling: its required
coincident pair exhausts the $A$-budget, so every other pair is
orthogonal.  This does not assert that arbitrary equality
configurations have this form; for instance, equality for three
distinct lines in a plane is also possible.
\end{document}
